\section{利用圆规画心脏线}
\markboth{71. 利用圆规画心脏线}{71. 利用圆规画心脏线}

从上节我们知道 $\varepsilon=1$ 时，以抛物线 $\rho=\dfrac{p}{1+\cos\theta}$ 的焦点为反演极，以 $R^2$ 为反演幂的反演图形是心脏线 $\rho=2r(1+\cos\theta)$，其中 $r=\dfrac{R^2}{2p}$。根据这个特点，可以利用圆规来画出心脏线。

\noindent
\begin{minipage}[t]{0.38\textwidth}
  \vspace{2pt}
  \centering
  \begin{tikzpicture}[scale=1.02, >=Stealth, line cap=round, line join=round]
    \def\p{1.15}
    \def\ym{0.42}
    \pgfmathsetmacro{\xa}{\p/2}
    \pgfmathsetmacro{\xp}{\p/2 - 2*\ym*\ym/\p}
    \pgfmathsetmacro{\yp}{2*\ym}

    % 极轴与直线 b
    \draw[->] (-0.65, 0) -- (1.15, 0) node[below, font=\scriptsize] {$x$};
    \draw (\xa, -1.35) -- (\xa, 1.35);
    \node[left=1pt, font=\scriptsize] at (\xa, -1.25) {$b$};

    % 抛物线 y^2 = 2p(p/2 - x) => x = p/2 - y^2/(2p)
    \draw[very thick, curveblue, domain=-1.18:1.18, samples=120, variable=\y]
      plot ({\p/2 - \y*\y/(2*\p)}, {\y});

    % 焦点 F, 顶点 A, 动点 M, 切点 P
    \coordinate (F) at (0,0);
    \coordinate (A) at (\xa, 0);
    \coordinate (M) at (\xa, \ym);
    \coordinate (P) at (\xp, \yp);

    % 线段 FM 与切线 MP
    \draw (F) -- (M);
    \draw ($(M)!-0.45!(P)$) -- ($(M)!1.65!(P)$);

    % 直角标记
    \draw ($(M)!0.13!(F)$) -- ($($(M)!0.13!(F)$) + ($(M)!0.13!(P)$) - (M)$) -- ($(M)!0.13!(P)$);

    % 节点标注
    \fill (F) circle (1.2pt) node[below left=0pt, font=\scriptsize] {$F$};
    \fill (A) circle (1.2pt) node[below right=0pt, font=\scriptsize] {$A$};
    \fill (M) circle (1.2pt) node[right=1pt, font=\scriptsize] {$M$};
    \fill (P) circle (1.2pt) node[above right=0pt, font=\scriptsize] {$P$};
  \end{tikzpicture}
  \captionof*{figure}{图\ 71.1}
\end{minipage}
\hfill
\begin{minipage}[t]{0.58\textwidth}
  \vspace{0pt}
  \hspace*{2em}在第 10 节中已经指出，抛物线关于焦点的垂足曲线是直线。由此可见，要作以定点 $F$ 为焦点的抛物线，只要如图 71.1 所示，作一条不通过 $F$ 的直线 $b$，在 $b$ 上任意取点 $M$，连 $FM$，作 $MP\perp FM$，则 $MP$ 是这抛物线的一条切线；这样作出若干条切线以后，再画它们的包络，就得到以 $F$ 为焦点的抛物线了。设这条抛物线的顶点是 $A$，则 $A$ 在 $b$ 上，并且 $FA\perp b$。
\end{minipage}

现在以 $F$ 为反演极，任意正数 $R^2$ 为反演幂，对图 71.1 中的整个图形作反演变换，结果得到图 71.2。

\medskip
\noindent
\begin{minipage}[t]{0.55\textwidth}
  \vspace{0pt}
  \hspace*{2em}在图 71.2 中，点 $A$、$M$、$P$ 分别变成 $B$、$N$、$Q$，直线 $AM$ 变成以 $FB$ 为直径的圆 $FBN$，直线 $MP$ 变成以 $FN$ 为直径的圆 $FNQ$，抛物线变成心脏线。因为抛物线与两直线 $MA$、$MP$ 都相切，所以得到的心脏线与两圆 $FBN$、$FNQ$ 都相切。但 $MA$ 就是定直线 $b$，所以它的反形 $FBN$ 是定圆，其圆心 $G$ 是极轴上的定点。设圆 $FNQ$ 的圆心是 $H$，连 $BN$、$GH$，则从 $FG=GB$、$FH=HN$，得 $GH\parallel BN$，所以
\end{minipage}
\hfill
\begin{minipage}[t]{0.42\textwidth}
  \vspace{2pt}
  \centering
  \begin{tikzpicture}[scale=0.88, >=Stealth, line cap=round, line join=round]
    % 心脏线参数 r=1.22 => B=(2.44,0), G=(1.22,0), 抛物线顶点 A=(0.28,0), p=2*xa=0.56, R^2 = r*p = 2*r*xa
    \def\r{1.22}
    \def\xa{0.28}
    \def\tzero{40}      % M 极角 40 deg, 使 F, A, M, H, N 及 P, Q 各自清晰分离

    \pgfmathsetmacro{\p}{2*\xa}
    \pgfmathsetmacro{\Rsq}{\r*\p}
    \pgfmathsetmacro{\fm}{\xa/cos(\tzero)}
    \pgfmathsetmacro{\fn}{2*\r*cos(\tzero)}
    \pgfmathsetmacro{\tq}{2*\tzero}
    \pgfmathsetmacro{\rq}{\r*(1+cos(\tq))}
    \pgfmathsetmacro{\rp}{\Rsq/\rq}

    \coordinate (F) at (0,0);
    \coordinate (A) at (\xa, 0);
    \coordinate (B) at (2*\r, 0);
    \coordinate (G) at (\r, 0);
    \coordinate (M) at ({\fm*cos(\tzero)}, {\fm*sin(\tzero)});
    \coordinate (N) at ({\fn*cos(\tzero)}, {\fn*sin(\tzero)});
    \coordinate (H) at ({0.5*\fn*cos(\tzero)}, {0.5*\fn*sin(\tzero)});
    \coordinate (Q) at ({\rq*cos(\tq)}, {\rq*sin(\tq)});
    \coordinate (P) at ({\rp*cos(\tq)}, {\rp*sin(\tq)});

    % 极轴
    \draw[->] (-0.82, 0) -- (2.72, 0) node[below, font=\scriptsize] {$x$};

    % 直线 b (虚线)
    \draw[densely dashed] (\xa, -1.55) -- (\xa, 1.55);
    \node[left=1pt, font=\scriptsize] at (\xa, -1.42) {$b$};

    % 抛物线 (虚线)
    \draw[densely dashed, domain=-1.12:1.12, samples=100, variable=\y]
      plot ({\xa - \y*\y/(2*\p)}, {\y});

    % 定圆 FBN (圆心 G, 半径 r)
    \draw (G) circle (\r);

    % 辅助定圆 (以 FG 为直径, 虚线)
    \draw[densely dashed] (0.5*\r, 0) circle (0.5*\r);

    % 圆 FNQ (圆心 H, 半径 fn/2)
    \draw (H) circle (0.5*\fn);

    % 心脏线 rho = r(1+cos t)
    \draw[very thick, curveblue, domain=0:360, samples=220, variable=\t]
      plot ({\r*(1+cos(\t))*cos(\t)}, {\r*(1+cos(\t))*sin(\t)});

    % 切线 MP (虚线) 与射线 FN, FQ, 弦 BN, 中位线 GH
    \draw[densely dashed] ($(M)!-0.45!(P)$) -- ($(M)!1.45!(P)$);
    \draw (F) -- (N) -- (B);
    \draw[densely dashed] (G) -- (H);
    \draw (F) -- (Q);

    % 标注点 (各自放置到清晰无压线方位)
    \fill (F) circle (1.2pt) node[below right=1pt, font=\scriptsize] {$F$};
    \fill (A) circle (1.2pt) node[below right=0pt, font=\scriptsize] {$A$};
    \fill (B) circle (1.2pt) node[below left=0pt, font=\scriptsize] {$B$};
    \fill (G) circle (1.2pt) node[below=1pt, font=\scriptsize] {$G$};
    \fill (H) circle (1.2pt);
    \node[font=\scriptsize] at ($(H)+(110:0.21)$) {$H$};
    \fill (M) circle (1.2pt) node[below right=0pt, font=\scriptsize] {$M$};
    \fill (N) circle (1.2pt) node[above right=0pt, font=\scriptsize] {$N$};
    \fill (P) circle (1.2pt) node[left=1pt, font=\scriptsize] {$P$};
    \fill (Q) circle (1.2pt) node[above left=0.5pt, font=\scriptsize] {$Q$};
  \end{tikzpicture}
  \captionof*{figure}{图\ 71.2}
\end{minipage}
\[
  \angle FHG = \angle FNB = 90^\circ.
\]

因而点 $H$ 在以 $FG$ 为直径的圆上。由于 $F$ 和 $G$ 都是定点，因而以 $FG$ 为直径的辅助圆是定圆。由此可得心脏线的画法：作以 $FG$ 为直径的圆，然后在圆周上任意取点 $H$，以点 $H$ 为圆心作通过 $F$ 的圆，这样作出若干个圆以后，再画它们的包络线，就得到心脏线了。
